\chapter{Selected Specialist Interview Excerpts}
\label{app:interviews}
\addcontentsline{toc}{chapter}{Appendix D: Selected Specialist Interview Excerpts}

\noindent
This appendix reproduces ten illustrative excerpts drawn from the
twenty-eight transcribed specialist interviews
(Section~\ref{sec:themes}, Chapter~\ref{ch:industry_reconfig}).  All
excerpts have been anonymised at the participant's request:
designations are reported in their generic form; specific film
titles, fee figures and platform-acquisition prices have been
redacted or aggregated where they would identify the speaker.  The
excerpts are presented in the order in which they were coded and
mapped onto the nine themes of Table~\ref{tab:themes}.

\section*{Excerpt 1 -- Senior producer (10--20 years' experience)}
\textit{Theme: Risk re-assessment of star-only logic}
\medskip

``I have been making films in this industry for almost twenty years.
The thing we believed for a very long time, which we admitted only
half-aloud, was that a top star is a kind of insurance.  You take a
script, you attach the right star, you sell the satellite, the
audio, the international, you get the theatre release, and even if
the film does not work creatively, the floor is high.  That logic
worked, more or less, until 2019.  After the lockdown, twice in
two years we saw mid-tier and even some upper-tier star films
collapse in the second weekend.  The reason was not that the
audience disappeared.  The audience came on day one because of the
star.  Then word-of-mouth on the script and the screenplay reached
WhatsApp groups by Saturday afternoon, and Sunday's collection was
forty per cent of Friday's.  That is a behaviour we did not see
before 2020 -- where the audience punishes a weak script very fast,
even with a star.  So now when we sit down with a writer and a
director, the first question is no longer `who is the star?'  It is
`what is the script?'  This is, for an old industry, a very large
shift.''

\section*{Excerpt 2 -- OTT acquisition head (5--10 years' experience)}
\textit{Theme: Hybrid window as new normal}
\medskip

``Three to four weeks is what works.  Less than three weeks and the
exhibitor walks out.  More than five weeks and we cannot recoup the
acquisition price because subscriber acquisition windows for the
title close out and Tamil-Nadu households who would have paid for
the streaming watch on YouTube clips and reviews instead.  At my
platform we have done about thirty-five Tamil titles in the last
two years and the median window we negotiated was twenty-six days.
The number is twenty-six days.  It varies a little.  Star films get
twenty-eight to thirty-five.  Mid-budget gets twenty-one to
twenty-eight.  The market has settled on this band almost without
anyone announcing it.  This is now the norm.  The producers are
comfortable with it because they get the OTT cheque earlier.  The
exhibitors are unhappy in principle but they have stopped fighting
it because, after the second weekend, occupancy collapses anyway in
most films, so the third week is not, in real terms, a high-revenue
week.''

\section*{Excerpt 3 -- Independent director (5--10 years' experience)}
\textit{Theme: Mid-budget renaissance / genre experimentation}
\medskip

``I made two films before 2020 -- both small, both struggled.  I am
making my third film now and the budget is around eleven crore.
The producer is comfortable with it; the script is genre material,
a slow-burn investigative thriller; we have a strong character lead
but no top star.  The OTT acquisition was conditionally agreed
before the shoot started.  Five years ago, this script would have
been impossible to make at this budget.  Either I would have had to
attach a star and inflate the budget to twenty-five crore, or the
producer would have asked me to write a comedy.  Now -- because the
audience reads investigative content on streaming, because the
acquisition price exists, because the exhibitor accepts a smaller
release -- the script is the entire pitch.  This is a genuine
opening for a generation of directors that did not exist before.''

\section*{Excerpt 4 -- Multiplex chain executive (10--20 years' experience)}
\textit{Theme: Modernisation of capacity / quality-weighted exhibition}
\medskip

``Our chain has not opened a new screen in Tamil Nadu since 2022.
What we have done instead is renovate.  Premium-format upgrades --
laser projection, Dolby Atmos, recliner seating -- have generated
better per-screen yield than fresh expansion.  The economics has
inverted: in 2018 a new screen in a Tier-2 town was a positive
investment; in 2024 the same screen in the same town pays back over
nine years instead of five.  Whereas a renovation in a high-traffic
multiplex pays back in eighteen months because the average ticket
moves from two hundred to three hundred and fifty rupees.  The
audience has not abandoned theatres.  It has become more
discriminating.  Quality of sound, quality of seating, quality of
food and beverage all matter to whether the audience comes out for
a film at all.''

\section*{Excerpt 5 -- Distributor (10--20 years' experience)}
\textit{Theme: Direct-to-OTT as portfolio play}
\medskip

``Five or six years ago, if a film went directly to OTT, it was a
failure.  We would say `the film could not get a release'.  That
narrative is gone.  About a third of the films I distributed in
2024 were structured as direct-to-streaming from inception.  These
are not failed theatrical films -- these are projects whose
economics make sense only on streaming, where the audience is
younger, more selective, and more genre-curious.  We negotiate the
acquisition price, we work with the platform on launch marketing,
and the producer's recovery is faster than a theatrical release
would have been.  The audience reaches it directly through the
platform's recommendations.  In some cases the audience is larger
than the audience we would have got with a small theatrical
release.  This is the new portfolio reality.''

\section*{Excerpt 6 -- Marketing head (10--20 years' experience)}
\textit{Theme: Marketing convergence on digital}
\medskip

``In 2018 we used to sit and plan the hoarding cities.  Five
hoardings in Madurai, three in Coimbatore, fifteen in Chennai, the
big one at LB Road, and so on.  The conversation was about real
estate.  Now the planning conversation is about which YouTuber, in
which language, with which audience, gets the embargo on the first
trailer.  We placed the audio launch teaser with about a dozen
mid-tier influencers two weeks before the launch and the first
twenty-four hours' reach exceeded the lifetime reach of any
hoarding campaign we did pre-2020.  The cost was a fraction.  This
is not a hypothesis any more, it is the operating reality.''

\section*{Excerpt 7 -- Music label executive (10--20 years' experience)}
\textit{Theme: Streaming-audio long-tail}
\medskip

``The audio launch is still magic for Tamil cinema.  Twenty
thousand people in a stadium, the composer on stage, the lead actor
in white and silver -- that has not changed.  What has changed is
where the money comes from.  In 2018 the gramophone-cassette-CD
revenue was already gone, but the satellite-rights and audio-rights
deal was a meaningful four to seven crore for an A-list film.  In
2024 the upfront audio rights have shrunk because the streaming
royalty is on a per-stream basis and the long-tail is what carries
the music.  A song that does forty million streams over twelve
months yields more royalty than a fixed audio-rights cheque, but it
arrives over a year, not on signature.  This has changed how we
finance the music itself, how we plan the release of the singles,
and how we structure the composer's deal.''

\section*{Excerpt 8 -- Exhibitor (single-screen, 20+ years' experience)}
\textit{Theme: Single-screen contraction and resilience}
\medskip

``Single-screen has lost a lot of ground.  In my circuit alone,
out of the twenty-three theatres I knew personally in 2019, eight
have closed permanently.  Four converted into commercial buildings;
two are running but on a cricket-and-events model; two are
permanently shut.  Of the remaining fifteen, eight are renovated
or partly renovated.  Seven are running on the old single-screen
model.  Will they survive ten years?  Honestly, I don't know.  The
star-fan-club opening day still gives us occupancy.  But the
mid-budget content film does not give us occupancy any more,
because that audience is now on streaming.  So we depend on three
or four big star releases a year to make our economics work, and
the rest of the year is barely covering rent and salaries.  The
government should help us with electricity tariffs and with GST,
but more important is what the producers do.  If they protect a
window of three weeks for theatrical, we can plan around it.  If
they don't, we cannot.''

\section*{Excerpt 9 -- Trade analyst (5--10 years' experience)}
\textit{Theme: Mid-tier star anxiety}
\medskip

``The thesis you should write into the academic literature is this:
the star system in Tamil cinema has not collapsed; it has
\emph{stratified}.  At the top, the four-five top names retain
opening-day pulling power and command fees that have corrected by
maybe ten to fifteen per cent.  Below them, the next ten or fifteen
A-list names have seen fees correct by twenty-five to thirty-five
per cent.  This is the band that is genuinely under pressure.  Why?
Because the audience that used to commit to an A-list film without
reading reviews now waits a day or two, sees the WhatsApp verdict,
and only then commits.  The premium an A-list star carries above a
director-led mid-budget film has shrunk.  This is the central
microeconomic fact of Tamil cinema today.  The character actors,
the supporting tier, are also under pressure but the reason there
is budget compression generally, not stardom dynamics
specifically.''

\section*{Excerpt 10 -- Pan-India producer (10--20 years' experience)}
\textit{Theme: Pan-India ambition}
\medskip

``Every project meeting in 2024 starts with a Pan-India question.
Can this be made to work in Telugu, Hindi, Malayalam, Kannada and
Tamil simultaneously?  Can we build a release in three to four
languages at once?  Should the lead actor have already proved
themselves outside Tamil before we attach them?  This was a
fringe question in 2019.  It is now the central question in the
greenlighting conversation.  The reasons are visible -- the
breakouts of \textit{Pushpa}, \textit{RRR}, \textit{KGF},
\textit{Kantara}, the success of multilingual streaming releases.
The pan-India construction does not work for every film, of course;
many small Tamil films are best made for Tamil only.  But for any
project above thirty crore the question must be asked, and the
answer increasingly drives the casting, the music, the marketing
and the windowing decisions.''

\bigskip
\noindent
\textit{End of Appendix D}
